\section{Moving logarithmic poles and cohomology comparisons}
\label{sec:logarithmic}

We now construct a single family through which the Higgs and Betti
correspondences can be compared.  The point carrying the full flag must move
over a compact curve.  Its motion changes the determinant, so the first step
is to compensate for that change by a line bundle.  The resulting comparison
will apply to the whole fixed-determinant cohomology.

\subsection{The determinant correction and the scalar correspondence}

Choose a point $x_0\in C$.  For the construction we first take
$L=\mathcal O_C(d x_0)$; we return to an arbitrary determinant at the end of
this section.  The line bundle $\mathcal O_C(d x)$ has a canonical
logarithmic connection $d_{d x}$: its canonical meromorphic section is
horizontal.  Its residue at $x$ is $-d$.  Multiplying a connection by
$\lambda$ produces a $\lambda$-connection, whose Leibniz rule is
$D(fs)=\lambda\,df\otimes s+fD(s)$.

Pull back the finite \emph{\'etale} map
$[n]:\operatorname{Pic}^0(C)\to\operatorname{Pic}^0(C)$ by
\[
 C\longrightarrow\operatorname{Pic}^0(C),\qquad
 x\longmapsto\mathcal O_C(d x-d x_0),
\]
and choose a connected component $S$ of this pullback.  Then $S$ is a smooth
projective connected curve, and its map $x:S\to C$ is finite \emph{\'etale}
and surjective.  Write $\Delta\subset C\times S$ for the graph of $x$.
A Poincar\'e bundle gives a line bundle $J$ on $C\times S$ such that
\begin{equation}
 J^{\otimes n}\simeq
 \mathcal O_{C\times S}\bigl(d\Delta-d(x_0\times S)\bigr)
       \otimes\operatorname{pr}_S^*B
 \label{log:root}
\end{equation}
for some line bundle $B$ on $S$.  The possible base factor is the only
ambiguity in this equality of families.

A line pulled back from $S$ has a canonical relative connection along $C$.
Equip the right side of \eqref{log:root} with the relative logarithmic
connection given by its canonical meromorphic section and this connection
on the base factor.  Dividing its local connection forms by $n$ defines
a unique compatible connection on $J$.  This division glues: the
transition equation for $J^{\otimes n}$ is $n$ times the equation for $J$.
Let $D_{J,\lambda}$ be the resulting relative $\lambda$-connection.  Its
residues are $-\lambda d/n$ along $\Delta$ and $+\lambda d/n$ along
$x_0\times S$.

Put
\[
 \mathfrak t_0=\{(u_1,\ldots,u_n)\in\mathbb A^n:
                         \textstyle\sum_i u_i=0\},\qquad
 \Lambda=\mathbb A^1_\lambda\times\mathfrak t_0.
\]
Let $\mathcal X\to\mathbb A^1_\lambda$ parametrize stable logarithmic
$\lambda$-connections on rank-$n$ bundles, with determinant
$(\mathcal O_C(d x_0),\lambda d_{d x_0})$ and residue
$-\lambda(d/n)I$ at $x_0$.  Stability here and below is slope stability for
subbundles preserved by the operator.  At $\lambda=0$ the scalar residue
is zero, so the fiber $X_0$ is the ordinary trace-free Higgs moduli space $X$.

Let $\mathcal Y\to S\times\Lambda$ parametrize the following data:
\begin{itemize}
 \item a stable logarithmic $\lambda$-connection $(E,D)$ with its only
 possible pole at $x(s)$ and determinant
 $(\mathcal O_C(d x(s)),\lambda d_{d x(s)})$;
 \item a full flag of $E_{x(s)}$, preserved by the residue, on whose
 ordered graded lines $Q_i$ the residue acts by $-\lambda d/n+u_i$.
\end{itemize}
The flag does not enter the stability condition.  We write $Y_{\lambda,u}$
for the fiber over $(\lambda,u)\in\Lambda$, \emph{retaining all of $S$}.
Thus $Y_{\lambda,u}$ is itself a family over the moving-point curve.

At $u=0$, let $\mathcal Z$ be the closed locus where the residue is scalar,
and let $\mathcal O$ be the same scalar-residue family with the flag
forgotten.  The correspondence over $\mathbb A^1_\lambda$ is shown
fiberwise in Figure~\ref{log:scalar-correspondence}.

\begin{figure}[htbp]
\centering
\begin{tikzcd}[column sep=large]
 X_\lambda & O_\lambda\simeq S\times X_\lambda
   \arrow[l,"\operatorname{pr}"'] & Z_\lambda
   \arrow[l,"\pi"'] \arrow[r,hook,"\iota"] & Y_{\lambda,0}
\end{tikzcd}
\caption{The scalar correspondence. All moving-pole spaces retain the
curve $S$; $\operatorname{pr}$ has fiber $S$, $\pi$ is a
flag projection of relative dimension $b$, and $\iota$ has codimension
$b$. The residue on $Y_{\lambda,0}$ may have a strictly flag-lowering
nilpotent part in addition to its prescribed scalar part; $Z_\lambda$
requires that nilpotent part to vanish. The derived morphisms and their cohomological and diagram-weight shifts
are recorded in Table~\ref{spec:shift-table}.}
\label{log:scalar-correspondence}
\end{figure}

Here $\pi$ forgets the flag and $\iota$ is the scalar-residue inclusion.
The first map comes from the following product identification.

\begin{lemma}
\label{log:scalar-product}
Tensoring by $(J,D_{J,\lambda})$ gives an algebraic isomorphism
$\mathcal O\simeq S\times\mathcal X$ over $S\times\mathbb A^1$.
At $\lambda=0$ it preserves the ordinary Hitchin coefficients.
\end{lemma}

\begin{proof}
By \eqref{log:root}, tensoring fixed-pole data by $J$ changes its
determinant to $\mathcal O_C(d x(s))$, up to a base line, and changes its
determinant operator to the prescribed one.  At $x_0$, the residue
$+\lambda d/n$ of $J$ cancels the residue $-\lambda d/n$ of the original
rank-$n$ connection.

This cancellation also holds where the moving and fixed sections meet.
In a coordinate $z$ with moving section $z=a(s)$, the scalar polar part
contributed by $J$ is
\[
 -\frac{\lambda d}{n}\frac{dz}{z-a(s)}
 +\frac{\lambda d}{n}\frac{dz}{z}.
\]
Adding the fixed-pole form leaves only the first term.  The identity is an
identity of relative meromorphic forms over the whole local base, including
$a(s)=0$.  Thus the tensor product has a single logarithmic pole at the
moving point.  The inverse construction tensors by $J^{-1}$.

Base-line twists do not change the induced map to coarse moduli, and
$\deg J_s=0$ preserves stability.  When $\lambda=0$, the Higgs field of
$J$ is zero, so the characteristic polynomial is unchanged.
\end{proof}

The relative quasiprojective moduli of logarithmic $\lambda$-connections
exist by the standard construction; we use the formulation in
\cite[Section~2.2]{dCFHLogHodge}.  Fixed determinant is obtained by taking
the inverse image of its rank-one section.  Adding a residue-invariant
full flag is projective over the unflagged moduli; its ordered residue
eigenvalues are then regular functions on the flag parameter space.
The prescribed-residue equations are closed.

We will use universal bundles only \emph{\'etale} locally.  Coprimality
makes semistable pairs stable, with scalar automorphisms.  Their stack is
a scalar gerbe over its coarse space, or a $\mu_n$-gerbe after choosing a
determinant isomorphism; compare
\cite[Lemma~5.3 and Remark~5.4]{dCFHLogHodge}.
The universal projective bundle, $\operatorname{End}E$, and the line
ratios $Q_i\otimes Q_j^{-1}$ have trivial scalar action and descend.
Individual $Q_i$ need not descend, and no construction below requires them
to do so.

\subsection{Smooth families and actual restriction isomorphisms}

Our next objective is to identify the cohomology of all fibers by
restriction from the total space.  Smoothness alone would not suffice,
because these moduli spaces are not proper.  The scaling action supplies
the additional compactness property that is needed.

\begin{proposition}
\label{log:families}
The morphism $\mathcal Y\to S\times\Lambda$ is smooth of relative
dimension $2R+2b$, and $\mathcal X\to\mathbb A^1$ is smooth of relative
dimension $2R$.  Their total spaces are semiprojective for scaling the
operators, which scales $\lambda$ and every $u_i$ with weight one.
The same holds for the inverse image in $\mathcal Y$ of every linear
subspace of $\Lambda$ through zero.  In
Figure~\ref{log:scalar-correspondence}, $\pi$ is a smooth proper flag
projection of relative dimension $b$, and $\iota$ is a regular embedding
of codimension $b$ with normal bundle $T_\pi^*$.
\end{proposition}

\begin{proof}
We first establish zero limits for scaling.  Given an unflagged object,
keep its bundle fixed and scale its operator.  This extends the orbit to
the stack of all logarithmic $\lambda$-connections, without imposing
semistability on the special fiber.  Semistable reduction
\cite[Proposition~4.48]{FHZ2025} replaces this DVR family by a semistable
one after a finite extension, without changing its generic fiber.
Its hypotheses hold for $\mathrm{GL}_n$ over $\mathbb C$.

The determinant of the replacement and the prescribed determinant agree
generically, hence agree over the DVR by separatedness of the rank-one
moduli space.  A generic flag extends because the flag variety of the
extended bundle at the pole is proper.  Residue preservation and the
\emph{ordered} diagonal residue equations are closed, and therefore
persist at the special fiber.  Coprimality makes the limiting semistable
pair stable.  The same reasoning applies to the scalar-residue condition.

A finite extension of the DVR is enough for a limit in the original
coarse space.  Indeed, embed this quasiprojective space in projective
space as a locally closed subvariety.  The original orbit has a unique
projective limit.  The limit after finite base change is that point and
belongs to the coarse moduli space, so the original orbit extends there.
This proves the required zero-limit property.

We next check smoothness over $\lambda=u=0$.  At such a point put
\[
 \mathcal F=\operatorname{ParEnd}_0(E),\qquad
 \mathcal G=\mathcal F^\vee\otimes K_C,
\]
where $\mathcal F$ is the sheaf of trace-free endomorphisms preserving
the flag at $x(s)$.  Trace pairing identifies $\mathcal G$ with the
trace-free logarithmic Higgs fields having strictly flag-lowering
residue.  The relative deformation complex is
\begin{equation}
 \mathcal F\xrightarrow{[\theta,-]}\mathcal G.
 \label{log:deformation}
\end{equation}

Here is the lifting argument, including the parameter directions.
For a small Artin extension with prescribed lifts of $s,\lambda,u$, the
bundle, determinant isomorphism and flag lift because the obstruction
lies in $H^2(C,\mathcal F)=0$.  In local frames respecting the lifted
flag, the operator lifts with its prescribed determinant and diagonal
residues: these are affine-linear conditions, and trace splits in
characteristic zero.  Its gluing obstruction lies in
$H^1(C,\mathcal G)$ tensored with the square-zero ideal.  Changing the
bundle and flag gluing changes this obstruction by the image of
$H^1(C,\mathcal F)$ under \eqref{log:deformation}.

This map on $H^1$ is surjective.  Its Serre dual is, up to sign, the map
on $H^0$ in \eqref{log:deformation}; its kernel consists of trace-free
automorphisms of the stable Higgs pair preserving the flag, hence is
zero.  Thus every such lifting problem has a solution.  Since
\[
 \chi(\mathcal F)=-(n^2-1)(g-1)-b=-R-b,
\]
the relative tangent dimension is $-2\chi(\mathcal F)=2R+2b$.
For the scalar unflagged family, replace $\mathcal F$ by
$\operatorname{End}_0(E)$ and obtain dimension $2R$ in the same way.
Finite scalar stabilizers act trivially on these deformations, so the
calculation also proves smoothness of the coarse spaces.

The nonsmooth locus of the morphism is closed and scaling invariant.
If it contained a point, it would contain that point's zero limit,
contradicting the calculation over $\lambda=u=0$.  This proves smoothness
everywhere.

A scaling-fixed point has $\lambda=u=0$.  Its logarithmic Hitchin
coefficients vanish, because their scaling weights are positive.
The logarithmic Hitchin map is proper; its base has coefficients in
$H^0(C,K_C(x(s))^{\otimes i})$ as in
\cite[Section~2.4, Equation~(4)]{dCFHLogHodge}.
The fixed locus is closed in its nilpotent fiber after adding flags
and imposing determinant and residue conditions.  Since flags and $S$
are proper, that fixed locus is proper.  Together with the zero limits,
this proves semiprojectivity.  The inverse image of a linear subspace
through zero is smooth by base change and inherits both properties.

Finally $\pi$ is the full flag bundle of the universal projective bundle.
On the $u=0$ family, subtracting the prescribed scalar from the residue
gives a section of the bundle of strictly flag-lowering matrices.
Its zero scheme is $\mathcal Z$.  The smoothness and dimensions just
proved show that it is a regular zero scheme of codimension $b$.
Trace pairing identifies this residue bundle along $\mathcal Z$ with
the cotangent bundle of the flag fibers, giving $N_\iota\simeq T_\pi^*$.
\end{proof}

\begin{corollary}
\label{log:restriction}
For every $(\lambda,u)\in\Lambda$, restriction induces an isomorphism
\begin{equation}
 \rho_{\lambda,u}:H^*(\mathcal Y,\mathbb Q)
       \xrightarrow{\ \sim\ }H^*(Y_{\lambda,u},\mathbb Q).
 \label{log:rho}
\end{equation}
Restriction to the inverse image of any line through zero in $\Lambda$
is also an isomorphism, as is restriction from that inverse image to any
of its fibers.  The analogous restriction isomorphisms hold for all
four families in Figure~\ref{log:scalar-correspondence} over the
$\lambda$-line.  They commute with pullbacks and with the proper Gysin
maps in that correspondence, including integration in the $S$ factor.
\end{corollary}

\begin{proof}
Recall that the core of a smooth semiprojective variety is the union
of the sets whose scaling orbits have limits as the parameter tends
to infinity.  Its inclusion induces a cohomology isomorphism
\cite[Theorem~1.3.1]{HRV2015}.  Because every coordinate on $\Lambda$
has positive weight, a point of $\mathcal Y$ with such a limit must
lie over zero.  Hence $\mathcal Y$ and its inverse image
$\mathcal Y_\ell$ of any line $\ell$ through zero have the same core.
Their restrictions to the core commute, so
$H^*(\mathcal Y)\to H^*(\mathcal Y_\ell)$ is an isomorphism.

The map $\mathcal Y_\ell\to\ell$ is smooth and surjective.  For
surjectivity, its image is a scaling-invariant open subset containing
zero, and therefore contains the whole affine line.  Apply
\cite[Corollary~1.3.3]{HRV2015}: for a smooth surjection from a smooth
semiprojective variety to the affine line, equivariant for a positive
base weight, restriction to each fiber is an isomorphism.  This proves
\eqref{log:rho} and its line version.  The same proof applies to
$\mathcal X$, $\mathcal O$ and $\mathcal Z$.

Pullback compatibility is functorial.  For Gysin maps, the families
are smooth over the line, the scalar embedding is transverse to fiber
restriction, and the flag and product projections are smooth proper.
The transverse base-change identities for Gysin maps give the claimed
compatibilities.
\end{proof}

\subsection{Riemann--Hilbert with a moving puncture}

Let $V$, $q$ and $Y^{\rm B}\to V$ be the ordered eigenvalue family from
Section~\ref{sec:betti}.  The exponent domain we need is
\begin{equation}
 U=\left\{u\in\mathfrak t_0^{\rm an}:
 t(u)=\bigl(\zeta e^{-2\pi\sqrt{-1}u_i}\bigr)_i\in V,\quad
 u_i-u_j\notin\mathbb Z\setminus\{0\}\right\}.
 \label{log:nonresonance}
\end{equation}
Notice that zero differences are allowed.  In particular $0\in U$,
which is essential for the scalar correspondence.

There is a flat bundle over $S$ with algebraic fiber family
$Y^{\rm B}\to V$.  To define it, use local trivializations of the
punctured surfaces $C\setminus\{x(s)\}$ and record a full flag of local
subsystems near the puncture.  We choose the peripheral-loop convention
in which its monodromy is the product of commutators used to define
$Y^{\rm B}$.  Changes of surface markings act on the handle matrices
by group words and on the peripheral flag by the corresponding
conjugating word.  They are algebraic maps over $V$.

These transition maps are locally constant and satisfy the cocycle
condition as algebraic maps on the coarse spaces.  To check this
precisely, let $\delta$ be the product of commutators in the free
fundamental group of the punctured surface.  A change of marking
gives an outer automorphism represented by $\varphi$, with
$\varphi(\delta)=w\delta w^{-1}$.  Word evaluation transforms a
representation to $\rho\varphi$ and its flag by $\rho(w)$.
Changing $\varphi$ by an inner automorphism gives simultaneous
conjugation.  Changing $w$ multiplies it on the right by an element
centralizing $\delta$.  The cyclically reduced word $\delta$ is not
a proper power: each positive handle generator occurs just once.
Its centralizer in the free group is therefore $\langle\delta\rangle$.
Thus this last ambiguity is a power of peripheral monodromy and
preserves every invariant flag, including those with repeated
eigenvalues or nonsemisimple monodromy.  The resulting maps compose
strictly after passage to the geometric quotient.  Consequently no
choice of eigenlines or higher homotopy of transitions is needed.
The scalar flag correspondence from
Section~\ref{sec:betti} gives a corresponding diagram of flat bundles
over $S$.

\begin{proposition}
\label{log:rh}
Over $U$, Riemann--Hilbert identifies $Y_{1,u}$ with the total space
over $S$ of this flat bundle at parameter $t(u)$.  These identifications
are homeomorphisms of families.  At $u=0$ they identify the scalar
correspondence of Figure~\ref{log:scalar-correspondence} at $\lambda=1$ with
the flat Betti scalar correspondence, compatibly with its pullbacks
and Gysin maps.
\end{proposition}

\begin{proof}
The local exponents of a connection are $r_i=-d/n+u_i$.
For the regular-singular correspondence and logarithmic lattices, including
Manin's extension theorem and Katz's nonresonance corollary, see
Deligne \cite[II, Proposition~5.4, Remark~5.5(ii),
Corollary~5.6, and Theorem~5.9]{Deligne1970}. We spell out the
nonresonant local construction to track the ordered flags and dependence
on the moving parameters.
In a disc coordinate $z$, write a logarithmic connection as
$d+A(z)\,dz/z$, with $A(0)=A_0$.  The regular-singular reduction
to $d+A_0\,dz/z$ by a gauge equal to the identity at zero is obtained
recursively using the operators $k+\operatorname{ad}(A_0)$, $k\geq1$.
Their eigenvalues are $k+r_i-r_j$, so they are invertible under
\eqref{log:nonresonance}, even when $A_0$ has repeated eigenvalues or
a nonzero nilpotent part.  On a small parameter neighborhood their
inverse norms are $O(1/k)$ for large $k$.  The usual power-series
majorant for the analytic coefficients therefore gives convergence
on a common smaller disc and continuous dependence on parameters.

In this constant-residue model a residue-invariant flag gives a
monodromy-invariant flag.  Conversely, the prescribed exponents select
one logarithm for each monodromy eigenvalue: equal exponential values
have equal selected exponents, by \eqref{log:nonresonance}.  Functional
calculus on each generalized eigenspace then gives a logarithm of
monodromy preserving its invariant flag.  This constructs the inverse
logarithmic extension with the specified ordered residues.  Coordinate
and logarithm changes in the constant-residue model act by functions
of the residue, and hence preserve the flag; the construction is
intrinsic.  In rank one the same construction gives exactly
$(\mathcal O_C(d x(s)),d_{d x(s)})$, so the determinant is the required
one.  The subproduct condition defining $V$ makes the representation
irreducible, and thus ensures stability.

Holonomy along local moving generators, together with this local
normal form, makes the forward maps continuous in the bundle and
in $s,u$.  Regard the Betti family as pulled back to $U$ via $t$;
there is no assertion that $t$ is globally injective.  The forward
map is a continuous bijection between manifolds of equal real
dimension, by Proposition~\ref{log:families} and the Betti
smoothness.  Invariance of domain gives its inverse continuity.
For a fixed point of $S$ it is holomorphic on the moduli fibers,
so it preserves their complex orientations.

At $u=0$, scalar residues correspond to scalar monodromy, and
forgetting or retaining the flag commutes with this construction.
The proper maps in the correspondence have their usual complex
orientations in the algebraic fibers and the same orientation on
$S$.  Consequently their topological Gysin maps agree under these
homeomorphisms.
\end{proof}

\begin{lemma}
\label{log:scalar-flat-product}
The product identification in Lemma~\ref{log:scalar-product}, at
$\lambda=1$, identifies the scalar unflagged Betti family with
$S\times X^{\rm B}$ as a flat bundle with algebraic fiber.  Its
projection to $X^{\rm B}$ corresponds to $\operatorname{pr}$.
\end{lemma}

\begin{proof}
Away from the fixed and moving poles, the nth power of the connection
on $J$ is the canonically trivial connection.  Its holonomies
therefore lie in $\mu_n$.  Locally in $S$, take a disc containing
$x_0$ and the moving point, and choose the handle generators outside
that disc.  Their line holonomies depend continuously on $s$ and
take values in the finite group $\mu_n$, so are locally constant.
Tensoring multiplies the handle matrices by these constant scalars,
an algebraic map on the representation variety.  The argument also
applies in collision neighborhoods.  Thus the algebraic product
identification of the scalar connection family is compatible with
the locally constant algebraic Betti transitions.
\end{proof}

We have now obtained the actual flat diagrams to which
Section~\ref{sec:diagram} will be applied.  In particular, when that
construction gives weight zero to the cochains of $S$, its product
and integration statements concern the same projection
$\operatorname{pr}$ as the logarithmic correspondence.

\subsection{The nonabelian-Hodge comparison on full cohomology}

The restriction isomorphisms for $\mathcal X$ and Riemann--Hilbert give
\begin{equation}
 H^*(X^{\rm B})\xrightarrow{\ \sim\ }H^*(X_1)
 \xleftarrow{\ \sim\ }H^*(\mathcal X)
 \xrightarrow{\ \sim\ }H^*(X).
 \label{log:comparison}
\end{equation}
We must check that this comparison transports the weight filtration
specified by nonabelian Hodge theory, including its nontrivial
$\Gamma$-character summands.  Agreement only after taking
$\Gamma$-invariants would not suffice.

\begin{lemma}
\label{log:finite-ambiguity}
Let a finite group $G$ act algebraically on a smooth connected complex
variety $B$, and let $q:B\to B/G$ be the quotient.  Suppose
$g:A\to B$ is a homeomorphism and $f:A\to B$ is continuous with
$qf=qg$.  Then $f=\gamma g$ for a single $\gamma\in G$.
\end{lemma}

\begin{proof}
Divide out the kernel of the action.  Smoothness and connectedness
make $B$ irreducible.  Removing the fixed loci of the nonidentity
elements gives a dense connected Zariski open $B^\circ$ on which
the action is free.  For $a\in g^{-1}(B^\circ)$ there is a unique
$\gamma(a)$ with $f(a)=\gamma(a)g(a)$.  The free finite quotient
is a covering there, so $\gamma(a)$ is locally constant.  It is
constant because this open set is connected.  Continuity extends
the equality over its closure, which is all of $A$.
\end{proof}

\begin{proposition}
\label{log:nah}
The comparison \eqref{log:comparison} transports the Deligne weight
filtration on the whole $H^*(X^{\rm B},\mathbb Q)$ in the same way as
the nonabelian-Hodge identification.  The conclusion is unchanged
when $\mathcal O_C(d x_0)$ is replaced by an arbitrary
$L\in\operatorname{Pic}^d(C)$ and the Higgs spaces are identified by
tensoring with a degree-zero nth root of their determinant ratio.
\end{proposition}

\begin{proof}
The full fixed-determinant nonabelian-Hodge diffeomorphism, including
its $\Gamma$-equivariance, is part of the twisted construction in
\cite[Sections~2 and~5]{HT2003}.  We compare its cohomology map with
the particular logarithmic family above.

For $(E,\theta)\in X$, choose a harmonic metric $H$ normalized to a
fixed metric on its determinant.  Let $\partial_H$ denote the
$(1,0)$ part of its Chern connection.  For real $0\leq\tau\leq1$
consider the holomorphic structure and logarithmic $\tau$-operator
\begin{equation}
 \bar\partial_E+\tau\theta^{\dagger_H},\qquad
 \tau\partial_H+\theta+
 \frac{\tau}{n}
 \bigl(d^{1,0}_{d x_0}-\partial_{\det H}\bigr)I.
 \label{log:harmonic-path}
\end{equation}
Its projectivization is the usual harmonic preferred section
\cite[Section~4]{Simpson1997}.  We verify the scalar adjustment.
The trace of $\theta$ and of $\theta^{\dagger_H}$ is zero, so the
determinant holomorphic structure stays fixed and the determinant
operator is exactly $\tau d_{d x_0}$.  The trace-free
holomorphicity equation follows from
$\bar\partial_E\theta=0$, $\partial_H\theta^{\dagger_H}=0$ and
the projective Hitchin equation.  The scalar equation is that of
the prescribed logarithmic determinant connection.  The only
nonregular term is the scalar logarithmic term at $x_0$.

For $\tau\ne0$, divide the operator by $\tau$.
Its puncture monodromy is $\zeta I$.  A proper invariant subspace
of rank $r$ would have $\zeta^r=1$ by taking determinants of the
commutator relation, contrary to $\gcd(n,d)=1$.  Thus these
connections are irreducible and stable.  At $\tau=0$ we recover
the original Higgs pair.  Hence \eqref{log:harmonic-path} lies in
$\mathcal X$ throughout.

This construction is a continuous homotopy of maps $X\to\mathcal X$.
One may work in local families of stable Higgs pairs.  The normalized
harmonic metrics vary smoothly: their trace-free metric linearization
is self-adjoint elliptic, and its kernel consists of infinitesimal
symmetries, which stability kills.  This is also the parameter
dependence used for preferred sections in
\cite[Section~4]{Simpson1997}.  To pass to the analytic topology of
coarse moduli, choose a sufficiently positive twist near any one
parameter so the underlying bundles are generated and have vanishing
$H^1$.  The kernels of the corresponding Dolbeault operators then
give continuously varying bases of sections.  Their evaluations
give continuous quotient presentations; the logarithmic operators,
viewed as morphisms from the first jet bundle with fixed symbol,
also have continuous coefficients in such presentations.  Passing
to the stable quotient gives the asserted continuity, independently
of the local choices.

At $\tau=0$ the homotopy is the inclusion $X\hookrightarrow\mathcal X$.
At $\tau=1$ it factors through $X_1$, and then through
Riemann--Hilbert to a map $F:X\to X^{\rm B}$.
The pullback $F^*$ is therefore precisely
\eqref{log:comparison}, because both restrictions from
$\mathcal X$ are isomorphisms.  Projectivizing $F$ gives the
usual projective nonabelian-Hodge map.

Let $\Phi:X\to X^{\rm B}$ be the fixed-determinant
nonabelian-Hodge diffeomorphism.  The maps $F$ and $\Phi$ agree
after projectivization.  Two determinant-one lifts of the same
projective representation differ on each handle generator by an
nth root of unity; thus their classes differ by the finite group
of determinant-preserving characters, identified with $\Gamma$.
The low-degree calculation of Section~\ref{sec:lefschetz} gives
$H^0(X^{\rm B})=\mathbb Q$, so $X^{\rm B}$ is connected; it is
smooth in the coprime setting.  Lemma~\ref{log:finite-ambiguity}
therefore gives $F=\gamma\Phi$ for one $\gamma\in\Gamma$.
Its action on $X^{\rm B}$ is algebraic and preserves Deligne
weights.  Consequently $F^*$ and $\Phi^*$ transport the weight
filtration identically on every cohomology group.  No averaging
or projection to invariants occurs.

Finally choose $M\in\operatorname{Pic}^0(C)$ with
$M^{\otimes n}\simeq L\otimes\mathcal O_C(-d x_0)$.
Tensoring by $M$ identifies the two Higgs moduli spaces
algebraically over the same Hitchin base.  Their projective
Higgs bundles, and therefore their projective harmonic
correspondences, are unchanged.  Applying the same finite-character
argument to the two fixed-determinant nonabelian-Hodge maps shows
that this identification also preserves the transported weight
filtration.  It preserves the perverse filtration because it is
an isomorphism over the Hitchin base.
\end{proof}
